\documentclass{beamer}
\usetheme{Madrid}
\usecolortheme{default}
\usepackage{amsmath, amssymb, graphicx, booktabs, array, multirow}
\usepackage[utf8]{inputenc}
\usepackage{hyperref}
\usepackage{tikz}

% Smaller font for more content per slide
\renewcommand{\tiny}{\fontsize{5pt}{6pt}\selectfont}
\renewcommand{\scriptsize}{\fontsize{6pt}{7pt}\selectfont}
\renewcommand{\footnotesize}{\fontsize{7pt}{8pt}\selectfont}
\renewcommand{\small}{\fontsize{8pt}{9pt}\selectfont}
\renewcommand{\normalsize}{\fontsize{9pt}{10pt}\selectfont}
\renewcommand{\large}{\fontsize{10pt}{11pt}\selectfont}

\title{CFA Level II: Capital Structure \& WACC}
\subtitle{Advanced WACC, CAPM, Fama-French \& Capital Structure Theory}
\author{CFA Level II Candidate Practice}
\date{2026}

\begin{document}
	
	\begin{frame}
		\titlepage
	\end{frame}
	
	\begin{frame}{Table of Contents}
		\footnotesize
		\tableofcontents
	\end{frame}
	
	% ============================================================================
	% SECTION 1: CAPITAL STRUCTURE THEORY - MM & TRADE-OFF
	% ============================================================================
	
	\section{Capital Structure Theory: MM \& Trade-Off}
	
	\begin{frame}{Q1: Modigliani-Miller Proposition I with Taxes}
		\scriptsize
		\textbf{Scenario:} Company A is an all-equity firm valued at \$500 million. The company is considering changing its capital structure by issuing \$200 million in debt at 6\% and using the proceeds to repurchase equity. The corporate tax rate is 30\%. The company's unlevered cost of equity is 12\%.
		
		\textbf{Question:} Using MM Proposition I with corporate taxes, calculate:
		1. The value of the levered firm
		2. The value of equity after the recapitalization
		3. The cost of equity after the recapitalization
		
		\vspace{2pt}
		\underline{\textbf{TRAP:}} MM Proposition I with taxes: V\_L = V\_U + tD. The value increases by the present value of the interest tax shield. Cost of equity increases with leverage: r\_e = r\_u + (r\_u - r\_d)(1-t)(D/E).
		
		\vspace{4pt}
		\underline{\textbf{Solution:}}
		
		\textbf{Step 1:} Value of levered firm
		V\_L = 500 + 0.30(200) = 500 + 60 = \$560M
		
		\textbf{Step 2:} Value of equity
		E = V\_L - D = 560 - 200 = \$360M
		
		\textbf{Step 3:} Cost of equity
		r\_e = 0.12 + (0.12 - 0.06)(0.70)(200/360)
		r\_e = 0.12 + 0.06(0.70)(0.5556) = 0.12 + 0.02333 = 0.14333 (14.33\%)
		
		\textbf{Step 4:} WACC
		WACC = (200/560)(0.06)(0.70) + (360/560)(0.1433)
		WACC = 0.015 + 0.0921 = 0.1071 (10.71\%)
		
		\textbf{Answer:} V\_L = \$560M; E = \$360M; r\_e = 14.33\%; WACC = 10.71\%.
	\end{frame}
	
	\begin{frame}{Q2: MM Proposition II with Bankruptcy Costs}
		\scriptsize
		\textbf{Scenario:} A company has an unlevered cost of equity of 10\%. The risk-free rate is 4\%, and the market risk premium is 6\%. The company is considering increasing its debt-to-equity ratio from 0.25 to 0.75. The cost of debt at the current leverage is 5\%, but increases to 7\% at the higher leverage. The tax rate is 35\%.
		
		\textbf{Question:}
		1. Calculate the cost of equity at both leverage levels using MM Proposition II
		2. Calculate WACC at both leverage levels
		3. Explain why the cost of equity increases more than proportionally at higher leverage
		
		\vspace{2pt}
		\underline{\textbf{TRAP:}} MM Proposition II: r\_e = r\_u + (r\_u - r\_d)(1-t)(D/E). The cost of equity increases with leverage, but the increase is moderated by the tax shield. At higher leverage, the cost of debt also increases.
		
		\vspace{4pt}
		\underline{\textbf{Solution:}}
		
		\textbf{Step 1:} Current leverage (D/E = 0.25)
		r\_e = 0.10 + (0.10 - 0.05)(0.65)(0.25) = 0.10 + 0.05(0.65)(0.25) = 0.10 + 0.008125 = 0.108125 (10.81\%)
		
		w\_d = 0.25/1.25 = 0.20, w\_e = 1.00/1.25 = 0.80
		
		WACC = 0.20(0.05)(0.65) + 0.80(0.108125) = 0.0065 + 0.0865 = 0.0930 (9.30\%)
		
		\textbf{Step 2:} Higher leverage (D/E = 0.75, r\_d = 7\%)
		r\_e = 0.10 + (0.10 - 0.07)(0.65)(0.75) = 0.10 + 0.03(0.65)(0.75) = 0.10 + 0.014625 = 0.114625 (11.46\%)
		
		w\_d = 0.75/1.75 = 0.4286, w\_e = 1.00/1.75 = 0.5714
		
		WACC = 0.4286(0.07)(0.65) + 0.5714(0.114625) = 0.0195 + 0.0655 = 0.0850 (8.50\%)
		
		\textbf{Step 3:} Explanation
		The cost of equity increases more at higher leverage because:
		1. Financial risk increases with leverage
		2. The cost of debt increases (5\% → 7\%) due to higher default risk
		3. The tax shield benefit is partially offset by higher bankruptcy risk
		
		\textbf{Answer:} Current: r\_e = 10.81\%, WACC = 9.30\%; Higher: r\_e = 11.46\%, WACC = 8.50\%.
	\end{frame}
	
	\begin{frame}{Q3: Optimal Capital Structure with Trade-Off Theory}
		\scriptsize
		\textbf{Scenario:} A company is evaluating its capital structure. The following data is available at different debt-to-equity ratios:
		
		\begin{tabular}{lrrrr}
			D/E Ratio & 0.00 & 0.50 & 1.00 & 1.50 \\
			Cost of Debt (pre-tax) & - & 5.0\% & 5.5\% & 6.5\% \\
			Unlevered Cost of Equity & 12\% & 12\% & 12\% & 12\% \\
			Tax Rate & 35\% & 35\% & 35\% & 35\% \\
		\end{tabular}
		
		The company has \$50 million in unlevered value.
		
		\textbf{Question:}
		1. Calculate the cost of equity, WACC, and firm value at each D/E ratio
		2. Identify the optimal capital structure
		3. Explain why increasing leverage beyond the optimal point destroys value
		
		\vspace{2pt}
		\underline{\textbf{TRAP:}} Optimal capital structure balances the tax shield benefit against the costs of financial distress. WACC is minimized at the optimal point.
		
		\vspace{4pt}
		\underline{\textbf{Solution:}}
	\end{frame}
	
	\begin{frame}{Q3: Solution}
		\scriptsize
		\textbf{Step 1:} Calculate for each D/E ratio
		
		\textbf{D/E = 0.00:}
		r\_e = 12\%
		WACC = 12\%
		V\_L = V\_U = \$50M
		
		\textbf{D/E = 0.50:}
		r\_e = 0.12 + (0.12 - 0.05)(0.65)(0.50) = 0.12 + 0.02275 = 0.14275 (14.275\%)
		
		w\_d = 0.50/1.50 = 0.3333, w\_e = 1.00/1.50 = 0.6667
		
		WACC = 0.3333(0.05)(0.65) + 0.6667(0.14275) = 0.01083 + 0.09517 = 0.1060 (10.60\%)
		
		V\_L = 50 + 0.35 × (50/1.50) = 50 + 11.67 = \$61.67M
		
		\textbf{D/E = 1.00:}
		r\_e = 0.12 + (0.12 - 0.055)(0.65)(1.00) = 0.12 + 0.04225 = 0.16225 (16.225\%)
		
		w\_d = 1.00/2.00 = 0.50, w\_e = 1.00/2.00 = 0.50
		
		WACC = 0.50(0.055)(0.65) + 0.50(0.16225) = 0.017875 + 0.081125 = 0.0990 (9.90\%)
		
		V\_L = 50 + 0.35 × (50/2.00) = 50 + 8.75 = \$58.75M
		
		\textbf{D/E = 1.50:}
		r\_e = 0.12 + (0.12 - 0.065)(0.65)(1.50) = 0.12 + 0.053625 = 0.173625 (17.3625\%)
		
		w\_d = 1.50/2.50 = 0.60, w\_e = 1.00/2.50 = 0.40
		
		WACC = 0.60(0.065)(0.65) + 0.40(0.173625) = 0.02535 + 0.06945 = 0.0948 (9.48\%)
		
		V\_L = 50 + 0.35 × (50/2.50) = 50 + 7.00 = \$57.00M
		
		\vspace{2pt}
		\textbf{Step 2:} Summary Table
		
		\begin{tabular}{lrrrr}
			D/E & r\_e & WACC & V\_L (M) \\
			0.00 & 12.00\% & 12.00\% & 50.00 \\
			0.50 & 14.28\% & 10.60\% & 61.67 \\
			1.00 & 16.23\% & 9.90\% & 58.75 \\
			1.50 & 17.36\% & 9.48\% & 57.00 \\
		\end{tabular}
		
		\vspace{2pt}
		\textbf{Step 3:} Optimal Capital Structure = D/E = 0.50 (highest firm value, lowest WACC)
		
		\textbf{Answer:} Optimal D/E = 0.50 with WACC = 10.60\% and V\_L = \$61.67M.
	\end{frame}
	
	\begin{frame}{Q4: Static Trade-Off Theory with Financial Distress Costs}
		\scriptsize
		\textbf{Scenario:} A company is analyzing its capital structure. The following data is available:
		
		\begin{tabular}{lrrrr}
			Debt (\$M) & 0 & 20 & 40 & 60 \\
			Interest Rate & - & 5.0\% & 5.5\% & 7.0\% \\
			Probability of Distress & 0\% & 5\% & 15\% & 30\% \\
			Expected Distress Cost (\% of Firm Value) & 0\% & 10\% & 20\% & 40\% \\
			Unlevered Firm Value & \$80M & - & - & - \\
			Tax Rate & 35\% & - & - & - \\
			Unlevered Cost of Equity & 11\% & - & - & - \\
		\end{tabular}
		
		\textbf{Question:}
		1. Calculate the levered firm value at each debt level considering both tax shields and distress costs
		2. What is the optimal debt level?
		3. Explain the implications of the static trade-off theory
		
		\vspace{2pt}
		\underline{\textbf{TRAP:}} V\_L = V\_U + PV(Tax Shield) - PV(Distress Costs). The optimal debt level balances the tax shield benefit against the expected distress costs.
		
		\vspace{4pt}
		\underline{\textbf{Solution:}}
	\end{frame}
	
	\begin{frame}{Q4: Solution}
		\scriptsize
		\textbf{Step 1:} Calculate levered value at each debt level
		
		\textbf{Debt = \$20M:}
		V\_tax = 80 + 0.35(20) = 87.00
		PV(Distress) = 0.05 × 0.10 × 87.00 = 0.435
		V\_L = 87.00 - 0.435 = \$86.565M
		
		r\_e = 0.11 + (0.11 - 0.05)(0.65)(20/66.565) = 0.11 + 0.06(0.65)(0.3005) = 0.11 + 0.01172 = 0.12172 (12.17\%)
		
		WACC = (20/86.565)(0.05)(0.65) + (66.565/86.565)(0.12172) = 0.0075 + 0.0936 = 0.1011 (10.11\%)
		
		\textbf{Debt = \$40M:}
		V\_tax = 80 + 0.35(40) = 94.00
		PV(Distress) = 0.15 × 0.20 × 94.00 = 2.82
		V\_L = 94.00 - 2.82 = \$91.18M
		
		r\_e = 0.11 + (0.11 - 0.055)(0.65)(40/51.18) = 0.11 + 0.055(0.65)(0.7816) = 0.11 + 0.02794 = 0.13794 (13.79\%)
		
		WACC = (40/91.18)(0.055)(0.65) + (51.18/91.18)(0.13794) = 0.0157 + 0.0774 = 0.0931 (9.31\%)
		
		\textbf{Debt = \$60M:}
		V\_tax = 80 + 0.35(60) = 101.00
		PV(Distress) = 0.30 × 0.40 × 101.00 = 12.12
		V\_L = 101.00 - 12.12 = \$88.88M
		
		r\_e = 0.11 + (0.11 - 0.07)(0.65)(60/28.88) = 0.11 + 0.04(0.65)(2.0776) = 0.11 + 0.0540 = 0.1640 (16.40\%)
		
		WACC = (60/88.88)(0.07)(0.65) + (28.88/88.88)(0.1640) = 0.0307 + 0.0533 = 0.0840 (8.40\%)
		
		\vspace{2pt}
		\textbf{Step 2:} Summary Table
		
		\begin{tabular}{lrrrr}
			Debt (M) & V\_L (M) & r\_e & WACC \\
			0 & 80.00 & 11.00\% & 11.00\% \\
			20 & 86.57 & 12.17\% & 10.11\% \\
			40 & 91.18 & 13.79\% & 9.31\% \\
			60 & 88.88 & 16.40\% & 8.40\% \\
		\end{tabular}
		
		\vspace{2pt}
		\textbf{Step 3:} Optimal Debt Level = \$40M (highest firm value)
		
		\textbf{Answer:} Optimal debt = \$40M, V\_L = \$91.18M, WACC = 9.31\%.
	\end{frame}
	
	\begin{frame}{Q5: Pecking Order Theory Application}
		\scriptsize
		\textbf{Scenario:} A company needs \$150 million in new capital for a positive NPV project. The company has the following characteristics:
		
		\begin{tabular}{lr}
			Current Debt & \$200M \\
			Current Equity & \$300M \\
			Net Income & \$50M \\
			Depreciation & \$30M \\
			Cash & \$25M \\
			Stock Price & \$20.00 \\
			Flotation Cost - Equity & 7\% \\
			Flotation Cost - Debt & 3\% \\
			Tax Rate & 35\% \\
			Cost of Debt (pre-tax) & 6\% \\
			Cost of Equity & 12\% \\
		\end{tabular}
		
		\textbf{Question:}
		1. According to the pecking order theory, what order of financing should the company follow?
		2. Calculate the effective cost of each financing source
		3. Why might the company prefer internal financing over external debt?
		
		\vspace{2pt}
		\underline{\textbf{TRAP:}} Pecking order theory: Internal financing (retained earnings) first, then debt, then equity as a last resort due to asymmetric information costs.
		
		\vspace{4pt}
		\underline{\textbf{Solution:}}
		
		\textbf{Step 1:} Available internal funds
		Retained Earnings = NI - Dividends = 50 - 20 = \$30M
		Cash available = 30 + 25 = \$55M (limited by cash balance)
		
		\textbf{Step 2:} Financing Sources
		1. Internal financing: \$55M (retained earnings + available cash)
		2. Debt financing: Remaining \$95M
		3. Equity financing: Not needed (last resort)
		
		\textbf{Step 3:} Effective cost of each source
		Internal financing cost = Opportunity cost = 12\% (what shareholders could earn elsewhere)
		
		After-tax cost of debt = 6\%(1-0.35) = 3.9\%
		Effective cost of debt with flotation = 3.9\%/(1-0.03) = 4.02\%
		
		Cost of equity with flotation = 12\%/(1-0.07) = 12.90\%
		
		\textbf{Step 4:} Pecking Order Preference
		1. Internal financing (cost = 12\% but no asymmetric information cost)
		2. Debt financing (cost = 4.02\%, lower information asymmetry)
		3. Equity financing (cost = 12.90\%, highest information asymmetry)
		
		\textbf{Answer:} Use internal funds first (\$55M), then debt (\$95M). Avoid equity issuance.
	\end{frame}
	
	% ============================================================================
	% SECTION 2: WACC CALCULATION
	% ============================================================================
	
	\section{WACC Calculation}
	
	\begin{frame}{Q6: WACC with Multiple Debt Issues}
		\scriptsize
		\textbf{Scenario:} A company has the following capital structure:
		
		\begin{tabular}{lrr}
			& Market Value & Cost \\
			Long-term Debt (Bonds) & \$150M & 6.5\% \\
			Bank Loan & \$80M & 5.0\% \\
			Preferred Stock & \$50M & 7.0\% \\
			Common Equity & \$320M & 11.5\% \\
		\end{tabular}
		
		The tax rate is 30\%. The company has 10 million shares outstanding at \$32 per share.
		
		\textbf{Question:}
		1. Calculate the market value weights of each component
		2. Calculate WACC
		3. If the company issues \$40M in new debt at 7.5\%, how would WACC change if the proceeds are used to repurchase shares at \$32?
		
		\vspace{2pt}
		\underline{\textbf{TRAP:}} Use market value weights, not book values. WACC = w\_d r\_d(1-t) + w\_p r\_p + w\_e r\_e.
		
		\vspace{4pt}
		\underline{\textbf{Solution:}}
	\end{frame}
	
	\begin{frame}{Q6: Solution}
		\scriptsize
		\textbf{Step 1:} Calculate market values and weights
		
		Total Market Value = 150 + 80 + 50 + 320 = \$600M
		
		Weights:
		w\_bonds = 150/600 = 0.25 (25\%)
		w\_bank = 80/600 = 0.1333 (13.33\%)
		w\_p = 50/600 = 0.0833 (8.33\%)
		w\_e = 320/600 = 0.5333 (53.33\%)
		
		\textbf{Step 2:} Combined debt cost and weight
		
		w\_d = 0.25 + 0.1333 = 0.3833 (38.33\%)
		r\_d = (150/230)(0.065) + (80/230)(0.05) = 0.6522(0.065) + 0.3478(0.05) = 0.0424 + 0.0174 = 0.0598 (5.98\%)
		
		\textbf{Step 3:} Calculate WACC
		
		WACC = 0.3833(0.0598)(0.70) + 0.0833(0.07) + 0.5333(0.115)
		WACC = 0.0161 + 0.0058 + 0.0613 = 0.0832 (8.32\%)
		
		\textbf{Step 4:} Impact of new debt issuance
		
		New debt = \$40M, new equity = 320 - 40 = \$280M
		Total = 150 + 80 + 50 + 280 = \$560M
		
		New weights:
		w\_d = 270/560 = 0.4821 (new total debt)
		w\_p = 50/560 = 0.0893
		w\_e = 280/560 = 0.5000
		
		New cost of debt (blended): r\_d = (150/270)(0.065) + (80/270)(0.05) + (40/270)(0.075)
		r\_d = 0.5556(0.065) + 0.2963(0.05) + 0.1481(0.075) = 0.0361 + 0.0148 + 0.0111 = 0.0620 (6.20\%)
		
		New WACC = 0.4821(0.0620)(0.70) + 0.0893(0.07) + 0.50(0.115)
		= 0.0209 + 0.0063 + 0.0575 = 0.0847 (8.47\%)
		
		\textbf{Answer:} Current WACC = 8.32\%; New WACC = 8.47\% (increases due to higher leverage).
	\end{frame}
	
	\begin{frame}{Q7: WACC with Convertible Debt and Warrants}
		\scriptsize
		\textbf{Scenario:} A company has the following capital structure:
		
		\begin{tabular}{lr}
			Convertible Debt (5\%, 10-year, convertible into 100 shares per \$1,000) & \$100M \\
			Straight Debt (6\%, 10-year) & \$200M \\
			Common Equity (20M shares at \$25) & \$500M \\
			Preferred Stock (6\%) & \$50M \\
			Tax Rate & 35\% \\
			Risk-Free Rate & 3.5\% \\
			Market Risk Premium & 6\% \\
			Beta (Levered) & 1.15 \\
		\end{tabular}
		
		\textbf{Question:}
		1. Calculate the cost of convertible debt considering its equity component
		2. Calculate WACC using the market value of the convertible debt
		3. Adjust the calculation if 30\% of the convertible debt is expected to be converted
		
		\vspace{2pt}
		\underline{\textbf{TRAP:}} Convertible debt has both debt and equity components. The cost should reflect the blended cost of the bond and conversion option.
		
		\vspace{4pt}
		\underline{\textbf{Solution:}}
	\end{frame}
	
	\begin{frame}{Q7: Solution}
		\scriptsize
		\textbf{Step 1:} Calculate cost of equity
		
		r\_e = R\_f + beta(ERP) = 0.035 + 1.15(0.06) = 0.035 + 0.069 = 0.104 (10.4\%)
		
		\textbf{Step 2:} Calculate cost of convertible debt
		
		Convertible value per bond = Straight bond value + Option value
		
		Straight bond value: N=10, I/Y=6, PMT=50, FV=1000 → PV = 92.64
		
		Option value = 100 × \$25 - 92.64 = 2,407.36
		
		Market value of convertible debt = \$100M (given)
		
		Implied cost of convertible debt: Solve 100 = 50 × PVIFA(r,10) + 1000 × PVIF(r,10) + Option Value
		
		Since the option value is embedded, we use the yield-to-maturity approach:
		r\_conv = 5.5\% (based on market price and conversion feature)
		
		\textbf{Step 3:} Calculate WACC with convertible debt as debt
		
		Total capital = 100 + 200 + 500 + 50 = \$850M
		
		w\_conv = 100/850 = 0.1176, w\_straight = 200/850 = 0.2353
		w\_e = 500/850 = 0.5882, w\_p = 50/850 = 0.0588
		
		WACC = 0.1176(0.055)(0.65) + 0.2353(0.06)(0.65) + 0.5882(0.104) + 0.0588(0.06)
		WACC = 0.0042 + 0.0092 + 0.0612 + 0.0035 = 0.0781 (7.81\%)
		
		\textbf{Step 4:} If 30\% of convertible debt is converted
		
		Converted portion: \$30M → equity; Remaining debt: \$70M
		
		New equity = 500 + 30 = \$530M; New debt = 70 + 200 = \$270M
		
		New weights: w\_d = 270/(270+530+50) = 270/850 = 0.3176; w\_e = 530/850 = 0.6235; w\_p = 0.0588
		
		WACC\_new = 0.3176(0.06)(0.65) + 0.6235(0.104) + 0.0588(0.06)
		= 0.0124 + 0.0648 + 0.0035 = 0.0807 (8.07\%)
		
		\textbf{Answer:} Current WACC = 7.81\%; After conversion = 8.07\%.
	\end{frame}
	
	\begin{frame}{Q8: Project-Specific WACC with Country Risk}
		\scriptsize
		\textbf{Scenario:} A US company is evaluating a project in Mexico. The company's domestic WACC is 8.5\%. The project has the following characteristics:
		
		\begin{tabular}{lr}
			Risk-Free Rate (US) & 3.0\% \\
			Risk-Free Rate (Mexico) & 6.0\% \\
			Market Risk Premium (US) & 5.0\% \\
			Market Risk Premium (Mexico) & 7.0\% \\
			Project Beta (Unlevered) & 1.20 \\
			Target D/E Ratio & 0.60 \\
			Cost of Debt (Mexico) & 8.0\% \\
			Tax Rate (Mexico) & 30\% \\
			Country Risk Premium (Mexico) & 2.5\% \\
			Inflation (US) & 2.0\% \\
			Inflation (Mexico) & 4.5\% \\
		\end{tabular}
		
		\textbf{Question:}
		1. Calculate the project's WACC using the country risk adjustment
		2. Should the project be evaluated using the domestic WACC? Why or why not?
		3. How does inflation affect the WACC calculation?
		
		\vspace{2pt}
		\underline{\textbf{TRAP:}} Project-specific WACC should reflect the project's risk and location, not the company's domestic WACC. Country risk must be incorporated.
		
		\vspace{4pt}
		\underline{\textbf{Solution:}}
	\end{frame}
	
	\begin{frame}{Q8: Solution}
		\scriptsize
		\textbf{Step 1:} Calculate levered beta considering country risk
		
		beta\_L = beta\_U[1 + (1-t)(D/E)] = 1.20[1 + (0.70)(0.60)] = 1.20(1.42) = 1.704
		
		\textbf{Step 2:} Calculate cost of equity with country risk
		
		r\_e = R\_f\^{US} + beta\_L(ERP\_US) + CRP
		r\_e = 0.03 + 1.704(0.05) + 0.025 = 0.03 + 0.0852 + 0.025 = 0.1402 (14.02\%)
		
		\textbf{Alternative:} Using local risk-free rate
		r\_e\^{local} = R\_f\^{Mexico} + beta\_L(ERP\_Mexico) = 0.06 + 1.704(0.07) = 0.06 + 0.1193 = 0.1793 (17.93\%)
		
		This should be converted to USD equivalent: r\_e\^{USD} = (1 + 0.1793)(1 + 0.02)/(1 + 0.045) - 1 = 0.1492 (14.92\%)
		
		\textbf{Step 3:} Calculate WACC
		
		WACC = w\_d r\_d(1-t) + w\_e r\_e
		w\_d = 0.60/1.60 = 0.375, w\_e = 1.00/1.60 = 0.625
		
		WACC = 0.375(0.08)(0.70) + 0.625(0.1402) = 0.0210 + 0.0876 = 0.1086 (10.86\%)
		
		\textbf{Step 4:} Why not use domestic WACC?
		- Domestic WACC (8.5\%) does not reflect:
		1. Higher country risk (political, economic, currency)
		2. Different inflation expectations
		3. Different tax rates
		4. Different market risk premiums
		
		\textbf{Step 5:} Inflation adjustment
		The real WACC = (1 + WACC nominal)/(1 + Inflation) - 1
		= (1.1086)/(1.045) - 1 = 0.0608 (6.08\%)
		
		\textbf{Answer:} Project WACC = 10.86\% (USD nominal). Use country-specific WACC, not domestic WACC.
	\end{frame}
	
	% ============================================================================
	% SECTION 3: FAMA-FRENCH MODELS
	% ============================================================================
	
	\section{Fama-French Models}
	
	\begin{frame}{Q9: Fama-French Three-Factor Model Application}
		\scriptsize
		\textbf{Scenario:} You are analyzing a stock with the following characteristics:
		
		\begin{tabular}{lr}
			Market Factor Beta & 1.20 \\
			SMB Beta & 0.45 \\
			HML Beta & -0.30 \\
			Risk-Free Rate & 3.0\% \\
			Expected Market Return & 9.0\% \\
			Expected Size Premium (SMB) & 2.5\% \\
			Expected Value Premium (HML) & 4.0\% \\
			Current Stock Price & \$45.00 \\
			Expected Dividend Next Year & \$2.00 \\
			Expected Dividend Growth Rate & 5\% \\
		\end{tabular}
		
		\textbf{Question:}
		1. Calculate the expected return using the Fama-French Three-Factor Model
		2. Compare with the CAPM expected return
		3. Is the stock undervalued or overvalued based on the Fama-French model?
		4. Why does the HML beta being negative affect the expected return?
		
		\vspace{2pt}
		\underline{\textbf{TRAP:}} The Fama-French model adds size and value factors. A negative HML beta means the stock behaves like a growth stock, reducing expected returns.
		
		\vspace{4pt}
		\underline{\textbf{Solution:}}
	\end{frame}
	
	\begin{frame}{Q9: Solution}
		\scriptsize
		\textbf{Step 1:} Fama-French Three-Factor Model
		
		r\_e = R\_f + beta\_M(ERP) + beta\_SMB(SMB) + beta\_HML(HML)
		
		ERP = 0.09 - 0.03 = 0.06 (6\%)
		
		r\_e = 0.03 + 1.20(0.06) + 0.45(0.025) - 0.30(0.04)
		r\_e = 0.03 + 0.072 + 0.01125 - 0.012 = 0.10125 (10.125\%)
		
		\vspace{2pt}
		\textbf{Step 2:} CAPM expected return
		
		r\_e\^{CAPM} = R\_f + beta(ERP) = 0.03 + 1.20(0.06) = 0.03 + 0.072 = 0.102 (10.2\%)
		
		\vspace{2pt}
		\textbf{Step 3:} Intrinsic value using Fama-French required return
		
		V\_0 = D\_1 / (r\_e - g) = 2.00 / (0.10125 - 0.05) = 2.00 / 0.05125 = \$39.02
		
		Current price = \$45.00 > \$39.02 → Stock is \textbf{overvalued}
		
		\vspace{2pt}
		\textbf{Step 4:} Why negative HML beta matters
		
		Negative HML beta means the stock has a negative exposure to the value premium:
		- When value stocks outperform growth stocks (HML positive), this stock underperforms
		- The stock behaves like a growth stock, which typically has lower expected returns
		- This is why the expected return is lower than CAPM would predict
		
		\vspace{2pt}
		\textbf{Step 5:} Comparison
		
		\begin{tabular}{lrr}
			& CAPM & FF3 \\
			Required Return & 10.20\% & 10.125\% \\
			Intrinsic Value & \$38.46 & \$39.02 \\
			Valuation & Overvalued & Overvalued \\
		\end{tabular}
		
		\textbf{Answer:} FF3 expected return = 10.125\%; Stock is overvalued at \$45.
	\end{frame}
	
	\begin{frame}{Q10: Fama-French Five-Factor Model Application}
		\scriptsize
		\textbf{Scenario:} A company has the following factor exposures:
		
		\begin{tabular}{lr}
			Market Beta & 0.95 \\
			SMB Beta & -0.15 \\
			HML Beta & 0.25 \\
			RMW Beta (Profitability) & 0.60 \\
			CMA Beta (Investment) & -0.40 \\
		\end{tabular}
		
		Risk-Free Rate = 2.5\%
		Expected Market Return = 8.0\%
		Expected SMB Premium = 2.0\%
		Expected HML Premium = 3.5\%
		Expected RMW Premium = 4.0\%
		Expected CMA Premium = 3.0\%
		
		The company has the following fundamentals:
		- Current Dividend = \$1.50
		- Expected Growth Rate = 4\%
		- Current Stock Price = \$28.00
		
		\textbf{Question:}
		1. Calculate the required return using the Fama-French Five-Factor Model
		2. Estimate the intrinsic value of the stock
		3. Compare with the Gordon growth model valuation
		4. Explain the implications of the negative CMA beta
		
		\vspace{2pt}
		\underline{\textbf{TRAP:}} The Five-Factor Model adds profitability (RMW) and investment (CMA) factors. Negative CMA means conservative investment policy.
		
		\vspace{4pt}
		\underline{\textbf{Solution:}}
	\end{frame}
	
	\begin{frame}{Q10: Solution}
		\scriptsize
		\textbf{Step 1:} Calculate ERP and individual factor premiums
		
		ERP = 0.08 - 0.025 = 0.055 (5.5\%)
		
		r\_e = R\_f + beta\_M(ERP) + beta\_SMB(SMB) + beta\_HML(HML) + beta\_RMW(RMW) + beta\_CMA(CMA)
		
		r\_e = 0.025 + 0.95(0.055) + (-0.15)(0.02) + 0.25(0.035) + 0.60(0.04) + (-0.40)(0.03)
		r\_e = 0.025 + 0.05225 - 0.003 + 0.00875 + 0.024 - 0.012
		r\_e = 0.0950 (9.50\%)
		
		\vspace{2pt}
		\textbf{Step 2:} Intrinsic value using Five-Factor required return
		
		V\_0 = D\_0(1+g) / (r\_e - g) = 1.50(1.04) / (0.095 - 0.04) = 1.56 / 0.055 = \$28.36
		
		Current price = \$28.00 → Stock is \textbf{undervalued} by 1.3\%
		
		\vspace{2pt}
		\textbf{Step 3:} Compare with Gordon growth model (using CAPM)
		
		CAPM: r\_e = 0.025 + 0.95(0.055) = 0.07725 (7.725\%)
		V\_0\^{GGM} = 1.56/(0.07725 - 0.04) = 1.56/0.03725 = \$41.88
		
		The GGM suggests the stock is significantly undervalued, but it ignores the size, value, profitability, and investment factors.
		
		\vspace{2pt}
		\textbf{Step 4:} Negative CMA beta implications
		
		Negative CMA beta (-0.40) means:
		- The company has a conservative investment policy (low capital expenditures relative to growth)
		- Conservative investment is typically associated with lower risk
		- This reduces the required return (since CMA premium is positive, negative beta × positive premium = negative contribution)
		
		\vspace{2pt}
		\textbf{Step 5:} Factor Contribution Analysis
		
		\begin{tabular}{lrr}
			Factor & Beta × Premium & Contribution \\
			Market & 0.95 × 5.5\% & 5.225\% \\
			SMB & -0.15 × 2.0\% & -0.300\% \\
			HML & 0.25 × 3.5\% & 0.875\% \\
			RMW & 0.60 × 4.0\% & 2.400\% \\
			CMA & -0.40 × 3.0\% & -1.200\% \\
			Total & - & 9.500\% \\
		\end{tabular}
		
		\textbf{Answer:} FF5 required return = 9.50\%; Intrinsic value = \$28.36; Slight undervaluation.
	\end{frame}
	
	\begin{frame}{Q11: Fama-French vs. CAPM in Valuation}
		\scriptsize
		\textbf{Scenario:} Two companies in the same industry:
		
		\begin{tabular}{lrr}
			& Company X & Company Y \\
			Market Beta & 1.10 & 1.05 \\
			SMB Beta & 0.50 & -0.20 \\
			HML Beta & 0.30 & 0.60 \\
			RMW Beta & 0.20 & 0.80 \\
			CMA Beta & 0.10 & -0.30 \\
			Current Dividend & \$2.00 & \$1.50 \\
			Growth Rate & 5\% & 4\% \\
			Stock Price & \$40.00 & \$30.00 \\
		\end{tabular}
		
		Risk-Free Rate = 3\%, ERP = 5.5\%, SMB = 2\%, HML = 3.5\%, RMW = 4\%, CMA = 3\%
		
		\textbf{Question:}
		1. Calculate required returns using CAPM and Fama-French Five-Factor
		2. Calculate intrinsic values using both models
		3. Which company is more attractively valued? Why?
		4. Explain the differences in factor exposures
		
		\vspace{2pt}
		\underline{\textbf{TRAP:}} The Fama-French model provides a more nuanced view of risk. Companies with different factor exposures may have different required returns even with similar betas.
		
		\vspace{4pt}
		\underline{\textbf{Solution:}}
	\end{frame}
	
	\begin{frame}{Q11: Solution}
		\scriptsize
		\textbf{Step 1:} CAPM Required Returns
		
		Company X: r\_e\^{CAPM} = 0.03 + 1.10(0.055) = 0.03 + 0.0605 = 0.0905 (9.05\%)
		
		Company Y: r\_e\^{CAPM} = 0.03 + 1.05(0.055) = 0.03 + 0.05775 = 0.08775 (8.775\%)
		
		\vspace{2pt}
		\textbf{Step 2:} Fama-French Required Returns
		
		Company X:
		r\_e\^{FF5} = 0.03 + 1.10(0.055) + 0.50(0.02) + 0.30(0.035) + 0.20(0.04) + 0.10(0.03)
		= 0.03 + 0.0605 + 0.01 + 0.0105 + 0.008 + 0.003 = 0.1220 (12.20\%)
		
		Company Y:
		r\_e\^{FF5} = 0.03 + 1.05(0.055) + (-0.20)(0.02) + 0.60(0.035) + 0.80(0.04) + (-0.30)(0.03)
		= 0.03 + 0.05775 - 0.004 + 0.021 + 0.032 - 0.009 = 0.12775 (12.775\%)
		
		\vspace{2pt}
		\textbf{Step 3:} Intrinsic Values
		
		\textbf{Company X:}
		D\_1 = 2.00(1.05) = 2.10
		CAPM: V\_0 = 2.10/(0.0905 - 0.05) = 2.10/0.0405 = \$51.85
		FF5: V\_0 = 2.10/(0.1220 - 0.05) = 2.10/0.072 = \$29.17
		
		\textbf{Company Y:}
		D\_1 = 1.50(1.04) = 1.56
		CAPM: V\_0 = 1.56/(0.08775 - 0.04) = 1.56/0.04775 = \$32.67
		FF5: V\_0 = 1.56/(0.12775 - 0.04) = 1.56/0.08775 = \$17.78
		
		\vspace{2pt}
		\textbf{Step 4:} Summary
		
		\begin{tabular}{lrrrr}
			& Price & CAPM Value & FF5 Value & CAPM vs FF5 \\
			Company X & \$40.00 & \$51.85 (UV) & \$29.17 (OV) & Large difference \\
			Company Y & \$30.00 & \$32.67 (UV) & \$17.78 (OV) & Large difference \\
		\end{tabular}
		
		\vspace{2pt}
		\textbf{Step 5:} Analysis
		
		The FF5 model suggests both companies are overvalued because:
		1. Company X has positive SMB (small cap premium) and positive CMA (conservative investment)
		2. Company Y has positive RMW (high profitability) but negative CMA
		3. The FF5 model captures additional risk factors that CAPM misses
		
		\textbf{Answer:} Both companies appear overvalued under FF5; CAPM would suggest undervaluation.
	\end{frame}
	
	% ============================================================================
	% SECTION 4: CAPITAL STRUCTURE DECISIONS
	% ============================================================================
	
	\section{Capital Structure Decisions}
	
	\begin{frame}{Q12: Leverage and Cost of Capital Trade-Off}
		\scriptsize
		\textbf{Scenario:} A company is considering its optimal capital structure. The following data is available:
		
		\begin{tabular}{lrrrr}
			Debt Ratio & 0\% & 20\% & 40\% & 60\% \\
			Cost of Debt & - & 5.0\% & 5.5\% & 7.0\% \\
			Beta (Levered) & 1.00 & 1.15 & 1.35 & 1.65 \\
		\end{tabular}
		
		Risk-Free Rate = 4\%, Market Risk Premium = 5\%, Tax Rate = 35\%
		
		\textbf{Question:}
		1. Calculate the cost of equity, WACC, and firm value at each debt ratio
		2. What is the optimal capital structure?
		3. How would the optimal capital structure change if the tax rate increases to 40\%?
		4. Explain the impact of the Modigliani-Miller propositions on this analysis
		
		\vspace{2pt}
		\underline{\textbf{TRAP:}} The optimal capital structure minimizes WACC and maximizes firm value. The cost of equity increases with leverage due to financial risk.
		
		\vspace{4pt}
		\underline{\textbf{Solution:}}
	\end{frame}
	
	\begin{frame}{Q12: Solution}
		\scriptsize
		\textbf{Step 1:} Calculate for each debt ratio
		
		\textbf{Debt Ratio = 0\%:}
		r\_e = 0.04 + 1.00(0.05) = 0.09 (9.0\%)
		WACC = 9.0\%
		V\_L = V\_U = \$100M (assuming)
		
		\textbf{Debt Ratio = 20\%:}
		r\_e = 0.04 + 1.15(0.05) = 0.04 + 0.0575 = 0.0975 (9.75\%)
		w\_d = 0.20, w\_e = 0.80
		WACC = 0.20(0.05)(0.65) + 0.80(0.0975) = 0.0065 + 0.078 = 0.0845 (8.45\%)
		
		V\_L = V\_U + tD = 100 + 0.35(20) = \$107.0M
		
		\textbf{Debt Ratio = 40\%:}
		r\_e = 0.04 + 1.35(0.05) = 0.04 + 0.0675 = 0.1075 (10.75\%)
		w\_d = 0.40, w\_e = 0.60
		WACC = 0.40(0.055)(0.65) + 0.60(0.1075) = 0.0143 + 0.0645 = 0.0788 (7.88\%)
		
		V\_L = 100 + 0.35(40) = \$114.0M
		
		\textbf{Debt Ratio = 60\%:}
		r\_e = 0.04 + 1.65(0.05) = 0.04 + 0.0825 = 0.1225 (12.25\%)
		w\_d = 0.60, w\_e = 0.40
		WACC = 0.60(0.07)(0.65) + 0.40(0.1225) = 0.0273 + 0.049 = 0.0763 (7.63\%)
		
		V\_L = 100 + 0.35(60) = \$121.0M
		
		\vspace{2pt}
		\textbf{Step 2:} Summary Table
		
		\begin{tabular}{lrrrr}
			Debt Ratio & r\_e & WACC & V\_L & Optimal? \\
			0\% & 9.00\% & 9.00\% & \$100M & No \\
			20\% & 9.75\% & 8.45\% & \$107M & No \\
			40\% & 10.75\% & 7.88\% & \$114M & No \\
			60\% & 12.25\% & 7.63\% & \$121M & Yes (Highest V\_L) \\
		\end{tabular}
		
		\vspace{2pt}
		\textbf{Step 3:} If tax rate increases to 40\%
		
		WACC = w\_d r\_d(0.60) + w\_e r\_e
		
		New WACC at 60\% debt: 0.60(0.07)(0.60) + 0.40(0.1225) = 0.0252 + 0.049 = 0.0742 (7.42\%)
		
		V\_L = 100 + 0.40(60) = \$124M
		
		Higher tax rate increases the value of the tax shield, making higher leverage more attractive.
		
		\textbf{Step 4:} MM Propositions Impact
		
		MM Proposition I (with taxes): V\_L = V\_U + tD (value increases with debt)
		MM Proposition II: r\_e = r\_u + (r\_u - r\_d)(1-t)(D/E) (cost of equity increases with leverage)
		
		The optimal capital structure balances the tax shield benefit against financial distress costs.
		
		\textbf{Answer:} Optimal debt ratio = 60\%; Higher tax rate increases optimal leverage.
	\end{frame}
	
	\begin{frame}{Q13: Leveraged Buyout (LBO) Capital Structure}
		\scriptsize
		\textbf{Scenario:} A private equity firm is considering an LBO of a target company. The target has:
		
		\begin{tabular}{lr}
			EBITDA (current) & \$200M \\
			EBIT & \$150M \\
			Debt (current) & \$300M \\
			Equity Value & \$500M \\
			Tax Rate & 35\% \\
			Cost of Equity (unlevered) & 12\% \\
		\end{tabular}
		
		The private equity firm plans to finance the acquisition with 70\% debt. The interest rate on the new debt is 7\%. The target is expected to grow at 5\% annually for 5 years, then stabilize at 3\%.
		
		\textbf{Question:}
		1. Calculate the new capital structure and WACC
		2. Estimate the enterprise value using the new capital structure
		3. What is the maximum price the private equity firm can pay?
		4. How does the increased leverage affect the cost of equity?
		
		\vspace{2pt}
		\underline{\textbf{TRAP:}} In an LBO, the capital structure changes significantly. The WACC should reflect the new leverage ratio and the associated costs.
		
		\vspace{4pt}
		\underline{\textbf{Solution:}}
	\end{frame}
	
	\begin{frame}{Q13: Solution}
		\scriptsize
		\textbf{Step 1:} Calculate new capital structure
		
		D/E = 0.70/0.30 = 2.333
		
		w\_d = 0.70, w\_e = 0.30
		
		\vspace{2pt}
		\textbf{Step 2:} Calculate levered beta (assuming unlevered beta = 0.85)
		
		beta\_L = beta\_U[1 + (1-t)(D/E)] = 0.85[1 + (0.65)(2.333)] = 0.85(1 + 1.516) = 0.85(2.516) = 2.139
		
		\vspace{2pt}
		\textbf{Step 3:} Calculate cost of equity
		
		r\_e = R\_f + beta\_L(ERP)
		Assuming R\_f = 4\%, ERP = 5\%:
		r\_e = 0.04 + 2.139(0.05) = 0.04 + 0.10695 = 0.14695 (14.695\%)
		
		\vspace{2pt}
		\textbf{Step 4:} Calculate WACC
		
		WACC = w\_d r\_d(1-t) + w\_e r\_e
		WACC = 0.70(0.07)(0.65) + 0.30(0.14695) = 0.03185 + 0.044085 = 0.075935 (7.594\%)
		
		\vspace{2pt}
		\textbf{Step 5:} Estimate enterprise value
		
		FCFF\_0 = EBITDA(1-t) - CapEx - WCInv = 200(0.65) - 50 - 30 = 130 - 80 = \$50M
		
		Using two-stage growth: 5\% for 5 years, then 3\% forever
		
		V\_firm = sum\_t=1 to 5 of FCFF\_t/(1.07594)\^t + FCFF\_6/((0.07594-0.03)(1.07594)\^5)
		
		FCFF\_1 = 50(1.05) = 52.5, FCFF\_2 = 55.1, FCFF\_3 = 57.9, FCFF\_4 = 60.8, FCFF\_5 = 63.8
		
		FCFF\_6 = 63.8(1.03) = 65.7
		
		V\_firm = 48.8 + 47.6 + 46.5 + 45.4 + 44.3 + 65.7/(0.04594)(0.07594)\^5
		= 232.6 + 65.7/(0.04594 × 0.680)
		= 232.6 + 65.7/0.0312 = 232.6 + 2,105.8 = \$2,338.4M
		
		\vspace{2pt}
		\textbf{Step 6:} Maximum price
		
		Maximum equity value = 2,338.4 - 0.70(2,338.4) = 701.5
		
		\textbf{Answer:} WACC = 7.594\%; Enterprise Value = \$2,338M; Maximum price = \$701.5M equity.
	\end{frame}
	
	\begin{frame}{Q14: Capital Structure Arbitrage}
		\scriptsize
		\textbf{Scenario:} A company has the following capital structure:
		
		\begin{tabular}{lr}
			Market Value of Debt & \$300M \\
			Market Value of Equity & \$500M \\
			Cost of Debt & 6\% \\
			Cost of Equity & 12\% \\
			Tax Rate & 35\% \\
			Unlevered Cost of Equity & 10\% \\
		\end{tabular}
		
		The company is considering increasing its debt by \$100M and using the proceeds to repurchase equity. After the recapitalization, the cost of debt is expected to increase to 6.5\%.
		
		\textbf{Question:}
		1. Calculate the current and new WACC
		2. Calculate the arbitrage profit opportunity if the company can increase leverage
		3. What is the optimal capital structure if financial distress costs are 5\% of firm value at the new leverage level?
		
		\vspace{2pt}
		\underline{\textbf{TRAP:}} Capital structure arbitrage exploits differences between actual and theoretical values. Financial distress costs reduce the benefit of additional leverage.
		
		\vspace{4pt}
		\underline{\textbf{Solution:}}
	\end{frame}
	
	\begin{frame}{Q14: Solution}
		\scriptsize
		\textbf{Step 1:} Current WACC
		
		w\_d = 300/800 = 0.375, w\_e = 500/800 = 0.625
		WACC = 0.375(0.06)(0.65) + 0.625(0.12) = 0.014625 + 0.075 = 0.089625 (8.9625\%)
		
		\vspace{2pt}
		\textbf{Step 2:} New capital structure
		
		New Debt = 400, New Equity = 400
		w\_d = 400/800 = 0.50, w\_e = 400/800 = 0.50
		
		\vspace{2pt}
		\textbf{Step 3:} New cost of equity using MM Proposition II
		
		r\_e\^{new} = r\_u + (r\_u - r\_d)(1-t)(D/E)\_new
		r\_e\^{new} = 0.10 + (0.10 - 0.065)(0.65)(400/400) = 0.10 + 0.02275 = 0.12275 (12.275\%)
		
		\vspace{2pt}
		\textbf{Step 4:} New WACC
		
		WACC\^{new} = 0.50(0.065)(0.65) + 0.50(0.12275) = 0.021125 + 0.061375 = 0.0825 (8.25\%)
		
		\vspace{2pt}
		\textbf{Step 5:} Arbitrage profit (without distress costs)
		
		Current value = 800, New value = Current + t(D\_new - D\_old) = 800 + 0.35(100) = \$835M
		
		Arbitrage profit = 835 - 800 = \$35M
		
		\vspace{2pt}
		\textbf{Step 6:} With financial distress costs
		
		Distress costs = 5\% of 835 = \$41.75M
		
		New value after distress = 835 - 41.75 = \$793.25M
		
		This is LESS than the current value of \$800M, so the recapitalization would destroy value.
		
		\vspace{2pt}
		\textbf{Step 7:} Optimal leverage
		
		Solve: V\_L = V\_U + tD - PV(Distress)
		
		At D = 400, V\_L = 800 + 35 - 41.75 = 793.25 < 800
		
		The optimal leverage is where tD = PV(Distress). This occurs at a lower debt level.
		
		\textbf{Answer:} Current WACC = 8.96\%; New WACC = 8.25\%; Arbitrage profit = \$35M (without distress); With distress costs, leverage should not be increased.
	\end{frame}
	
	% ============================================================================
	% SECTION 5: ADVANCED WACC APPLICATIONS
	% ============================================================================
	
	\section{Advanced WACC Applications}
	
	\begin{frame}{Q15: WACC with Inflation and Real Options}
		\scriptsize
		\textbf{Scenario:} A company is evaluating a project with the following characteristics:
		
		\begin{tabular}{lr}
			Project Life & 5 years \\
			Initial Investment & \$100M \\
			Expected Annual Cash Flow (Year 1) & \$30M \\
			Real Growth Rate & 3\% \\
			Nominal Growth Rate & 5.5\% \\
			Risk-Free Rate (nominal) & 4\% \\
			Market Risk Premium & 5\% \\
			Project Beta & 1.20 \\
			Tax Rate & 35\% \\
			Target D/E Ratio & 0.50 \\
			Cost of Debt (nominal) & 6.5\% \\
		\end{tabular}
		
		\textbf{Question:}
		1. Calculate the nominal and real WACC
		2. Evaluate the project using both nominal and real cash flows
		3. How would the presence of a real option (option to expand) affect the valuation?
		4. What is the project's NPV?
		
		\vspace{2pt}
		\underline{\textbf{TRAP:}} Nominal WACC should be used with nominal cash flows; real WACC with real cash flows. The Fisher equation links nominal and real rates.
		
		\vspace{4pt}
		\underline{\textbf{Solution:}}
	\end{frame}
	
	\begin{frame}{Q15: Solution}
		\scriptsize
		\textbf{Step 1:} Calculate cost of equity
		
		r\_e = R\_f + beta(ERP) = 0.04 + 1.20(0.05) = 0.04 + 0.06 = 0.10 (10\%)
		
		\vspace{2pt}
		\textbf{Step 2:} Calculate nominal WACC
		
		w\_d = 0.50/1.50 = 0.3333, w\_e = 1.00/1.50 = 0.6667
		
		WACC\_nominal = 0.3333(0.065)(0.65) + 0.6667(0.10)
		= 0.01408 + 0.06667 = 0.08075 (8.075\%)
		
		\vspace{2pt}
		\textbf{Step 3:} Calculate real WACC using Fisher equation
		
		WACC\_real = (1 + WACC\_nominal)/(1 + Inflation) - 1
		
		Inflation = 5.5\% - 3\% = 2.5\%
		
		WACC\_real = 1.08075/1.025 - 1 = 1.05439 - 1 = 0.05439 (5.439\%)
		
		\vspace{2pt}
		\textbf{Step 4:} Nominal cash flows
		
		Year 1: 30.00
		Year 2: 30(1.055) = 31.65
		Year 3: 31.65(1.055) = 33.39
		Year 4: 33.39(1.055) = 35.23
		Year 5: 35.23(1.055) = 37.16
		
		NPV\_nominal = -100 + 30/1.08075 + 31.65/1.08075\^2 + 33.39/1.08075\^3 + 35.23/1.08075\^4 + 37.16/1.08075\^5
		= -100 + 27.76 + 27.09 + 26.44 + 25.81 + 25.19 = 32.29 (\$M)
		
		\vspace{2pt}
		\textbf{Step 5:} Real cash flows
		
		Real cash flows = Nominal cash flows / (1.025)\^t
		
		Year 1: 29.27, Year 2: 30.12, Year 3: 31.00, Year 4: 31.90, Year 5: 32.82
		
		NPV\_real = -100 + 29.27/1.05439 + 30.12/1.05439\^2 + 31.00/1.05439\^3 + 31.90/1.05439\^4 + 32.82/1.05439\^5
		= -100 + 27.76 + 27.09 + 26.44 + 25.81 + 25.19 = 32.29 (\$M)
		
		\vspace{2pt}
		\textbf{Step 6:} Real option adjustment
		
		Option to expand adds value. If expansion value = \$15M, then:
		NPV\_with\_option = 32.29 + 15 = \$47.29M
		
		\textbf{Answer:} Nominal WACC = 8.075\%; Real WACC = 5.439\%; NPV = \$32.29M; With real option = \$47.29M.
	\end{frame}
	
	\begin{frame}{Q16: WACC and Business Risk vs. Financial Risk}
		\scriptsize
		\textbf{Scenario:} A company has the following characteristics:
		
		\begin{tabular}{lr}
			Unlevered Beta & 0.90 \\
			Risk-Free Rate & 3.5\% \\
			Market Risk Premium & 5.5\% \\
			D/E Ratio & 0.40 \\
			Tax Rate & 30\% \\
			Cost of Debt (pre-tax) & 5.5\% \\
			Current WACC & 9.2\% \\
		\end{tabular}
		
		The company is considering a project with higher business risk. The project has:
		- Unlevered Beta = 1.20
		- Project D/E Ratio = 0.60
		- Project Cost of Debt = 6.5\%
		
		\textbf{Question:}
		1. Decompose the current WACC into business risk and financial risk components
		2. Calculate the project's cost of equity and WACC
		3. Should the company use the project WACC or the corporate WACC?
		
		\vspace{2pt}
		\underline{\textbf{TRAP:}} Business risk is reflected in the unlevered beta; financial risk is reflected in the leverage. Project-specific WACC should reflect the project's risk profile.
		
		\vspace{4pt}
		\underline{\textbf{Solution:}}
	\end{frame}
	
	\begin{frame}{Q16: Solution}
		\scriptsize
		\textbf{Step 1:} Current cost of equity
		
		beta\_L = beta\_U[1 + (1-t)(D/E)] = 0.90[1 + (0.70)(0.40)] = 0.90(1.28) = 1.152
		
		r\_e = 0.035 + 1.152(0.055) = 0.035 + 0.06336 = 0.09836 (9.836\%)
		
		\vspace{2pt}
		\textbf{Step 2:} Verify WACC
		
		w\_d = 0.40/1.40 = 0.2857, w\_e = 1.00/1.40 = 0.7143
		
		WACC = 0.2857(0.055)(0.70) + 0.7143(0.09836) = 0.0110 + 0.0703 = 0.0813 (8.13\%)
		
		The given WACC is 9.2\%, indicating there may be additional company-specific risk.
		
		\vspace{2pt}
		\textbf{Step 3:} Decompose WACC
		
		Business risk component (unlevered): r\_u = R\_f + beta\_U(ERP) = 0.035 + 0.90(0.055) = 0.0845 (8.45\%)
		
		Financial risk component: r\_e - r\_u = 0.09836 - 0.0845 = 0.01386 (1.386\%)
		
		Tax shield benefit: w\_d r\_d t = 0.2857(0.055)(0.30) = 0.00471 (0.471\%)
		
		\vspace{2pt}
		\textbf{Step 4:} Project cost of equity
		
		beta\_L\^{proj} = 1.20[1 + (0.70)(0.60)] = 1.20(1.42) = 1.704
		
		r\_e\^{proj} = 0.035 + 1.704(0.055) = 0.035 + 0.09372 = 0.12872 (12.872\%)
		
		w\_d = 0.60/1.60 = 0.375, w\_e = 1.00/1.60 = 0.625
		
		WACC\^{proj} = 0.375(0.065)(0.70) + 0.625(0.12872) = 0.01706 + 0.08045 = 0.09751 (9.751\%)
		
		\vspace{2pt}
		\textbf{Step 5:} Comparison and recommendation
		
		\begin{tabular}{lrr}
			& Corporate & Project \\
			WACC & 8.13\% - 9.2\% & 9.751\% \\
		\end{tabular}
		
		The project has higher business risk (beta 1.20 vs 0.90) and higher financial risk (D/E 0.60 vs 0.40), so it requires a higher WACC.
		
		\textbf{Answer:} Use project WACC = 9.751\% (not corporate WACC). Business risk component = 8.45\%, financial risk = 1.386\%.
	\end{frame}
	
	\begin{frame}{Q17: WACC and Marginal Cost of Capital}
		\scriptsize
		\textbf{Scenario:} A company has the following capital structure and costs:
		
		\begin{tabular}{lr}
			Current Debt & \$200M at 5.5\% \\
			Current Equity & \$300M at 11\% \\
			Tax Rate & 35\% \\
		\end{tabular}
		
		The company needs to raise \$100M for a new project. Financing options:
		1. Additional debt at 6.0\%
		2. Additional equity at 12.5\%
		3. Preferred stock at 7.5\%
		
		The company has retained earnings of \$30M available for investment.
		
		\textbf{Question:}
		1. Calculate the current WACC
		2. Calculate the marginal cost of capital (MCC) for each financing option
		3. Determine the optimal financing mix if the project has an IRR of 10\%
		4. What is the breakpoint in the marginal cost of capital?
		
		\vspace{2pt}
		\underline{\textbf{TRAP:}} The marginal cost of capital is the cost of raising additional capital. The MCC schedule can have breakpoints where the cost changes.
		
		\vspace{4pt}
		\underline{\textbf{Solution:}}
	\end{frame}
	
	\begin{frame}{Q17: Solution}
		\scriptsize
		\textbf{Step 1:} Current WACC
		
		w\_d = 200/500 = 0.40, w\_e = 300/500 = 0.60
		WACC = 0.40(0.055)(0.65) + 0.60(0.11) = 0.0143 + 0.066 = 0.0803 (8.03\%)
		
		\vspace{2pt}
		\textbf{Step 2:} Marginal cost of each option
		
		\textbf{Option A - Additional Debt:}
		After-tax cost = 6.0\%(1-0.35) = 3.90\%
		
		\textbf{Option B - Additional Equity:}
		Cost = 12.5\%
		
		\textbf{Option C - Preferred Stock:}
		Cost = 7.5\%
		
		\textbf{Option D - Retained Earnings:}
		Cost = 11\% (opportunity cost)
		
		\vspace{2pt}
		\textbf{Step 3:} MCC Schedule with weights
		
		If using retained earnings first (\$30M):
		MCC\_RE = 0.40(0.055)(0.65) + 0.60(0.11) = 0.0803 (8.03\%)
		
		If using debt for remaining \$70M:
		MCC\_debt = 0.40(0.06)(0.65) + 0.60(0.11) = 0.0156 + 0.066 = 0.0816 (8.16\%)
		
		If using preferred stock:
		MCC\_preferred = 0.40(0.055)(0.65) + 0.60(0.11) = 0.0803 (with preferred, weights may change)
		
		\vspace{2pt}
		\textbf{Step 4:} Breakpoint Analysis
		
		Breakpoint = Amount of cheaper capital / Weight of that capital in capital structure
		
		Breakpoint = 30/0.60 = \$50M
		
		At \$50M, the MCC increases because retained earnings are exhausted.
		
		\vspace{2pt}
		\textbf{Step 5:} Project Acceptance Decision
		
		Project IRR = 10\% > MCC = 8.16\%, so the project should be accepted.
		
		\vspace{2pt}
		\textbf{Step 6:} Optimal Financing Mix
		
		1. Use \$30M retained earnings (cost = 11\% pre-tax)
		2. Use \$70M debt (cost = 3.90\% after-tax)
		3. Weighted MCC = 0.30(0.11) + 0.70(0.039) = 0.033 + 0.0273 = 0.0603 (6.03\%) after-tax
		
		\textbf{Answer:} Current WACC = 8.03\%; MCC increases after \$50M; Project should be accepted (IRR = 10\% > MCC).
	\end{frame}
	
\end{document}